\chapter*{Phụ Lục II. Mười Hai Điển Tích Để Soi Lại Mình}
\addcontentsline{toc}{chapter}{Phụ Lục II. Mười Hai Điển Tích Để Soi Lại Mình}
\markboth{Mười Hai Điển Tích}{Mười Hai Điển Tích}

\section*{1. Gia Cát Lượng và chữ Kiệm}
\begin{truyen}{Kiệm Dĩ Dưỡng Đức}{Frugality Nurtures Virtue}
Gia Cát Lượng xem tiết kiệm là cách nuôi dưỡng đức hạnh. Điều ấy nhắc ta rằng tiền bạc và nhân cách không hề tách rời nhau.\
\textit{(Zhuge Liang saw frugality as a way to nourish virtue. This reminds us that money and character are not separate matters.)}
\end{truyen}

\section*{2. Quản Ninh không rời chỗ ngồi}
\begin{truyen}{Không Đứng Dậy Trước Sự Hào Nhoáng}{Staying Seated Before Splendor}
Quản Ninh không chạy theo xe ngựa sang trọng. Bài học là giữ đường của mình trước ánh nhìn bên ngoài.\
\textit{(Guan Ning did not run after the splendid carriage. The lesson is to keep one’s path despite outward glitter.)}
\end{truyen}

\section*{3. Tăng Tử giữ lời với con}
\begin{truyen}{Trẻ Con Học Từ Điều Người Lớn Làm}{Children Learn from Adult Action}
Tăng Tử giữ chữ tín bằng hành động. Trong chuyện tiền bạc, con cái cũng học từ nếp sống của cha mẹ như thế.\
\textit{(Zengzi kept his promise through action. In matters of money, children likewise learn from the habits of parents.)}
\end{truyen}

\section*{4. Người nước Tề nhổ mạ}
\begin{truyen}{Sốt Ruột Làm Hỏng Việc}{Impatience Ruins Growth}
Sự sốt ruột khiến người ta làm hỏng cả ruộng mạ. Trong tiêu tiền, sốt ruột cũng làm hỏng khả năng tích lũy.\
\textit{(Impatience ruined the farmer’s whole field. In spending, impatience can ruin one’s capacity to accumulate.)}
\end{truyen}

\section*{5. Tri túc giả phú}
\begin{truyen}{Người Biết Đủ Mới Giàu}{Contentment Is Real Wealth}
Cổ nhân xem biết đủ là giàu. Khi lòng thôi thiếu mãi, đồng tiền mới ở lại lâu hơn.\
\textit{(The ancients saw contentment as true wealth. When the heart stops feeling endlessly deprived, money stays longer.)}
\end{truyen}

\section*{6. Liệu cơm gắp mắm}
\begin{truyen}{Sống Theo Sức Mình}{Living According to One’s Means}
Đây là trí khôn mộc mạc nhưng bền vững: đừng để thể diện buộc mình phải sống vượt quá khả năng.\
\textit{(This is simple but durable wisdom: do not let pride force you to live beyond your means.)}
\end{truyen}

\section*{7. Tích cốc phòng cơ}
\begin{truyen}{Chuẩn Bị Trước Khi Thiếu}{Prepare Before Lack Arrives}
Dành trước cho ngày khó là thái độ trưởng thành, không phải tâm thế sợ hãi.\
\textit{(Setting aside for hard days is maturity, not fearfulness.)}
\end{truyen}

\section*{8. Một hạt gạo một sợi chỉ}
\begin{truyen}{Nhớ Nguồn Gốc Của Vật}{Remember the Origin of Things}
Ai nhớ công sức phía sau từng món đồ sẽ khó lãng phí hơn.\
\textit{(Whoever remembers the labor behind each object will waste less.)}
\end{truyen}

\section*{9. Lão Tử và bước chân đầu tiên}
\begin{truyen}{Bắt Đầu Nhỏ Nhưng Bắt Đầu Thật}{Begin Small but Begin for Real}
Thói quen tiết kiệm không chờ điều kiện hoàn hảo. Nó chờ một bước chân thật sự.\
\textit{(The habit of saving does not wait for perfect conditions. It waits for one real step.)}
\end{truyen}

\section*{10. Gia thư Tăng Quốc Phiên}
\begin{truyen}{Đi Lên Nhưng Đừng Xa Hoa Quá Sớm}{Rise Without Early Luxury}
Người đi lên dễ hỏng ở chỗ vừa khá hơn một chút đã quên nề nếp cũ.\
\textit{(Those who are rising often fail by forgetting discipline as soon as life improves a little.)}
\end{truyen}

\section*{11. Chữ Nhẫn trong đời sống thường ngày}
\begin{truyen}{Nhịn Không Phải Là Thiếu Bản Lĩnh}{Restraint Is Not Weakness}
Biết nhịn một món đồ, một cuộc vui hay một cơn sĩ diện là bản lĩnh của người nghĩ xa.\
\textit{(To refrain from a purchase, a pleasure, or a moment of pride is the strength of one who thinks far ahead.)}
\end{truyen}

\section*{12. Hậu sinh khả úy, nhưng phải học từ gốc}
\begin{truyen}{Muốn Con Vững, Người Lớn Phải Vững Trước}{If Children Are to Stand Firm, Adults Must Stand First}
Người trẻ có thể giỏi hơn cha mẹ, nhưng nền nếp ban đầu vẫn được trao từ mái nhà.\
\textit{(The young may surpass their parents, but their earliest discipline is still handed down from home.)}
\end{truyen}

\begin{baihoc}
Điển tích không để trưng bày học rộng. Điển tích có ích nhất khi soi được vào thói quen rất nhỏ của hôm nay.
\textit{(Classical anecdotes are not for displaying scholarship. They are most useful when they illuminate the smallest habits of today.)}
\end{baihoc}

\begin{ghinhoanh}
Ancient wisdom becomes useful

only when it changes present habits.
\end{ghinhoanh}

\ngancach